% !TEX program = xelatex

%%%%%%%%%%%%%%%%%%%%%%%%%%%%%%%%
%
% This is a template for a hebrew document in latex. 
% It is to be used to typeset lovely documents in the holey language
%
% version 1.0
%
% License: CC BY 4.0 
%
% Author: Elad Denenberg
%%%%%%%%%%%%%%%%%%%%%%%%%%%%%%%%%
%
% NOTICE:
%   * The document currently is built so that the defualt language is Hebrew. 
%     Therefore all titles and text are Right to Left (RTL) by default
%   * This setting is good if you're writing mostly in Hebrew and have some English text or math
%   * If you're writing mainly in English and need some Hebrew it is advised to use English as the default language
%%%%%%%%%%%%%%%%%%%%%%%%%%%%%%%%%%

\documentclass{article} 
%\usepackage[utf8]{inputenc}

\usepackage{polyglossia} %This is the package that gives access to fonts.
\setdefaultlanguage{hebrew}  %this makes the titles RTL, switch these if you write mostly in English
\setotherlanguage{english}  
\newfontfamily\hebrewfont[Script=Hebrew]{David CLM} %This font will work in Overleaf, Linux and Mac installs, if you're running windows you'd need to pick the right ones
\newfontfamily\hebrewfontsf[Script=Hebrew]{Miriam CLM}

 %notice the \textenglish command here. This is because otherwise the RTL would screw up the parenthesis 
\title{\textenglish{Hebrew Document in LaTeX (actually XeLaTeX)} עברית בלטעך}
\author{אלעד דננברג}
\date{יג באב תשפ"ב \textenglish{August 2022}}

\begin{document}

\maketitle

\section{חלק: מבוא} %Hebrew title, on the right

לורם איפסום דולור סיט אמט, קונסקטורר אדיפיסינג אלית קולורס מונפרד אדנדום סילקוף, מרגשי ומרגשח. עמחליף נולום ארווס סאפיאן - פוסיליס קוויס, אקווזמן קוואזי במר מודוף. אודיפו בלאסטיק מונופץ קליר, בנפת נפקט למסון בלרק - וענוף לפרומי בלוף קינץ תתיח לרעח. לת צשחמי קולהע צופעט למרקוח איבן איף, ברומץ כלרשט מיחוצים. קלאצי סחטיר בלובק. תצטנפל בלינדו למרקל אס לכימפו, דול, צוט ומעיוט - לפתיעם ברשג - ולתיעם גדדיש. קוויז דומור ליאמום בלינך רוגצה. לפמעט

\subsection{תת חלק}

הועניב היושבב שערש שמחויט - שלושע ותלברו חשלו שעותלשך וחאית נובש ערששף. זותה מנק הבקיץ אפאח דלאמת יבש, כאנה ניצאחו נמרגי שהכים תוק, הדש שנרא התידם הכייר וק.

\begin{equation}\label{eq:visc}
\tau = \mu \frac{\partial u}{\partial y}
\end{equation}

ליבם סולגק. בראיט ולחת צורק מונחף, בגורמי מגמש. תרבנך וסתעד לכנו סתשם השמה - לתכי מורגם בורק? לתיג ישבעס.

\subsubsection{תת תת חלק}

פיסקא שמערבת עברית ואנגלית ניתן לכתוב כך:\textenglish{Some English text}. ניתן לכתוב גם פיסקאות שלומות באנגלית, אנחנו נרצה שהן יהיו משמאל לימין כי הן באנגלית. שימו לב לפיסקא הבאה, היא מתחילה מצד שמאל. 

\begin{english}
    Lorem ipsum dolor sit amet, consectetur adipiscing elit, sed do eiusmod tempor incididunt ut labore et dolore magna aliqua. Ut enim ad minim veniam, quis nostrud exercitation ullamco laboris nisi ut aliquip ex ea commodo consequat. Duis aute irure dolor in reprehenderit in voluptate velit esse cillum dolore eu fugiat nulla pariatur. Excepteur sint occaecat cupidatat non proident, sunt in culpa qui officia deserunt mollit anim id est laborum
\end{english}

והפיסקה הזו חזרה להיות מצד ימין שוב. בהצלחה בכתיבת טקסט בעברית !!

\end{document}

